\chapter{Distinctions and Distinguishability Loss}
\label{ch:distinctions}

\begin{quotation}
\noindent\textit{Synopsis.} This chapter establishes the formal vocabulary
used throughout the book. Drawing on Distinguishability Geometry, it treats
information as a structured space of distinctions and describes renderings
as projections that preserve some distinctions while collapsing others. The
chapter's central move is the distinction between benign underdetermination
and corruption: not all loss is harmful, and corruption arises not from
loss alone but from \emph{undeclared} loss that obscures what has been
discarded. This reframes information quality around disclosure rather than
preservation alone.
\end{quotation}

\section{Distinguishability Geometry, Imported}

Distinguishability Geometry supplies this monograph's most basic object: a
space of distinctions $(X,\sim_X)$, together with an equivalence structure
recording what is and is not told apart within it. Nothing in this section
modifies that prior formalism; it is specialized here to the case where
$X$ is the distinction-space of an object $e$ under transformation.

\begin{remark}[A Standing Simplification]
\label{rem:finite-distinctions}
Throughout this monograph, distinction spaces are treated as finite or
countable, and the cardinality-based measures used later
(\cref{def:fidelity} in particular) are stated only for that case. This is
a simplification made for tractability, not a claim that every domain of
interest is discrete in the relevant sense; a painting's causal-history
space, for instance, is arguably continuous rather than merely large. Where
this matters, it is flagged again locally (\cref{ch:painting}); it is not
resolved here.
\end{remark}

\section{Renderings as Projections}

\begin{definition}[Rendering as Projection]
\label{def:rendering}
A rendering $r$ of an object $e$ is a map $(X,\sim_X) \to (Y,\sim_Y)$
under which some distinctions separated in $X$ become identified in $Y$.
\end{definition}

This definition is stated at the level of Distinguishability Geometry's
own primitives: $X$ and $Y$ are sets of possible states, and $\sim_X$,
$\sim_Y$ record which pairs of states are, and are not, told apart within
each. This is the correct level of generality for the formalism, but it is
not the level at which the rest of this monograph actually talks. From
Chapter~\ref{ch:admissibility} onward, and especially throughout
Chapter~\ref{ch:budget-relay}, this book speaks of \emph{a} distinction
$d$ -- ``distinction $d_3$,'' ``the repair cost of $d$,'' ``$d$ is dropped
at stage $i$'' -- as though a distinction were a single named thing, not a
pair of states. Reconciling these two ways of talking, rather than
quietly switching between them, is this chapter's principal piece of
technical work.

\subsection{From Pairs to Named Distinctions}

\begin{definition}[Discarded Pair]
\label{def:discarded-pair}
Given a rendering $r: (X,\sim_X) \to (Y,\sim_Y)$ realized by an underlying
map $\varphi: X \to Y$, an unordered pair $\{x, x'\} \subseteq X$ is a
\emph{discarded pair} of $r$ if $x \not\sim_X x'$ (the pair is told apart
in the source) but $\varphi(x) \sim_Y \varphi(x')$ (the pair becomes
indistinguishable under $r$).
\end{definition}

\cref{def:discarded-pair} is the pair-level primitive Distinguishability
Geometry actually provides, and it is the correct formal object underlying
everything that follows. It is, however, an awkward unit to reason about
directly: nobody asks whether ``the pair \{recording with emphasis on the
third syllable, recording without it\}'' was preserved; they ask whether
\emph{emphasis} was preserved. The bridge from pairs to named features is
given by treating a feature as the thing that distinguishes a minimal pair.

\begin{definition}[Minimal Pair for a Distinction]
\label{def:minimal-pair}
A named distinction $d$ is represented, relative to a source $(X,\sim_X)$,
by a minimal pair $\{x_d, x_{\neg d}\} \subseteq X$: two states differing
exactly in whether the feature named by $d$ holds, and told apart in the
source ($x_d \not\sim_X x_{\neg d}$).
\end{definition}

\begin{definition}[Discarded Distinction Set, Revised]
\label{def:discard-set}
Given a rendering $r$ and a fixed assignment of minimal pairs to named
distinctions (\cref{def:minimal-pair}), the discarded distinction set
$\discard{r}$ is the set of named distinctions $d$ whose minimal pair
$\{x_d, x_{\neg d}\}$ is a discarded pair of $r$ (\cref{def:discarded-pair}):
\[
\discard{r} \;=\; \{\, d \;:\; \{x_d, x_{\neg d}\} \text{ is discarded by } r \,\}.
\]
\end{definition}

\cref{def:discard-set} supersedes an earlier, looser formulation of the
same definition, which described $\discard{r}$ directly as ``the set of
distinction-pairs separated in the source but identified in $r$'' and was
consequently ambiguous between the pair-level and feature-level readings
used elsewhere in the book. Every later chapter's usage --
\cref{def:benign-loss} and \cref{def:corrupting-loss} below,
Chapter~\ref{ch:repair-economics}'s treatment of $\repcost{d}$ for a named
$d$, and the entire apparatus of Chapter~\ref{ch:budget-relay} -- presumes
the feature-level reading given here, and \cref{def:discard-set} is stated
so as to make that presumption explicit and correct rather than
convenient and unexamined.

\begin{remark}[When the Bridge Is Clean]
\label{rem:bridge-clean}
The identification in \cref{def:discard-set} between a discarded pair and
a discarded feature is unproblematic exactly when $X$ factors, at least
locally, as a product of independent feature-dimensions -- so that varying
$d$ alone, holding everything else fixed, is a meaningful operation on $X$
and produces a genuine minimal pair rather than an artifact of an
arbitrary choice of coordinates. This holds cleanly for the toy and
worked examples used throughout this monograph (a transcript's lexical
content and prosodic emphasis are reasonably treated as independent
dimensions of an audio recording) and is assumed as a standing
simplification rather than argued for in general. Domains where features
are not independently variable in this sense -- where, for instance, no
change to one property is possible without simultaneously changing
another -- are not treated here.
\end{remark}

\section{Underdetermination Versus Corruption}

The distinction developed in this section is used, in one form or
another, in every subsequent chapter of the monograph, and is the closest
thing this book has to a single organizing move.

\begin{definition}[Benign Loss]
\label{def:benign-loss}
A discarded distinction $d \in \discard{r}$ is benign relative to a task
$\tau$ if $d$ is not required to perform $\tau$ correctly.
\end{definition}

\begin{definition}[Corrupting Loss]
\label{def:corrupting-loss}
A discarded distinction $d \in \discard{r}$ is corrupting relative to
$\tau$ if $\tau$ requires $d$ and the rendering's presentation obscures
that $d$ was discarded, as opposed to openly declaring the gap.
\end{definition}

\begin{proposition}[Corruption Requires Undeclared Loss]
\label{prop:corruption-undeclared}
Let $r$ be a rendering with discard set $\discard{r}$. If every
distinction in $\discard{r}$ is declared -- that is, made available to the
consumer as a known gap rather than presented as preserved -- then loss
alone does not imply corruption. Corruption requires both (i) the loss of
a task-relevant distinction, and (ii) failure to disclose that loss.
\end{proposition}

\begin{proofsketch}
Corruption, per \cref{def:corrupting-loss}, is task-relative by
construction. If $d\in\discard{r}$ is declared, a consumer pursuing a task
requiring $d$ can compensate for its absence, decline to rely on $r$ for
that task, or seek an alternative rendering; the distinction remains
unavailable, but the consumer's belief-state about $r$'s sufficiency is
not falsified by $r$ itself. Undeclared loss is therefore necessary for
corruption in the sense of \cref{def:corrupting-loss}, though not
sufficient, since a task not requiring $d$ is unaffected regardless of
disclosure.
\end{proofsketch}

\cref{def:benign-loss} and \cref{def:corrupting-loss} are stated here in
terms of ``required to perform $\tau$ correctly'' without yet saying
precisely what correctness of performance means, since that requires the
apparatus of admissible continuation sets developed in
Chapter~\ref{ch:admissibility}. \cref{prop:benign-preserves-sufficiency}
there closes the loop, restating both definitions in the sharper,
task-relative vocabulary Chapter~\ref{ch:admissibility} provides.

\section{A Toy Example}

\begin{example}
\label{ex:jpeg}
Lossy image compression discards high-frequency detail below a
quantization threshold. When the compression ratio and quantization
parameters are attached to the file as metadata, the loss is declared: a
consumer needing pixel-exact reproduction can detect, before relying on
the image, that it will not get it. When the same file is relabeled as a
lossless master, an identical discard set becomes corrupting under
\cref{def:corrupting-loss}, though nothing about the image data itself has
changed. In the vocabulary of \cref{def:minimal-pair}, each discarded
high-frequency component corresponds to a minimal pair of source images
differing only in that component, rendered indistinguishable once
quantized; the metadata either does or does not disclose that this
minimal pair has been collapsed. The example is deliberately kept toy at
this stage; it is developed properly, without simplification, in the
media-pipeline case of Chapter~\ref{ch:pipelines}.
\end{example}
